% Section: schemas

\documentclass[../../../include/open-logic-section]{subfiles}

\begin{document}

\olfileid[fr]{nml}{syn}{sch}

\olsection{Schémas et validité}

\begin{defn}
  Un \emph{schéma} est un ensemble de !!{formula}s formé de
  toutes les instances, et d'elles seules, obtenues par
  substitution à partir d'une !!{formula} modale~$!C$, soit
  \[
  \Setabs{!B}{\lexists[!D_1], \dots, \lexists[!D_n] \left(!B =
  \SSubst{!C}{\subst{!D_1}{p_1}, \dots, \subst{!D_n}{p_n}} \right) }.
  \]
  La !!{formula}~$!C$ est appelée \emph{formule caractéristique}
  du schéma; elle est unique à un renommage des !!{propositional
  variable}s près. Une !!{formula}~$!A$ est une \emph{instance}
  d'un schéma si elle appartient à cet ensemble.
\end{defn}

Il est commode de désigner un schéma par l'expression métalinguistique
obtenue en remplaçant les composantes atomiques de~$!C$ par
« $!A$ », « $!B$ »,~\dots. Ainsi, « $!A$ »,
« $!A \lif \Box!A$ » et « $!A \lif (!B \lif !A)$ »
désignent des schémas. Leurs !!{formula}s caractéristiques sont
respectivement $p$, $p \lif \Box p$ et $p \lif (q \lif p)$.
Le schéma « $!A$ » désigne l'ensemble de \emph{toutes} les
!!{formula}s.

\begin{defn}
  Un schéma est \emph{vrai} dans un modèle si et seulement si
  toutes ses instances y sont vraies; il est \emph{valide} si et
  seulement s'il est vrai dans tous les modèles.
\end{defn}

\begin{prop}\ollabel{prop:Kvalid}
  Le schéma \Ax{K} suivant est valide :
  \begin{equation*}
  \Box(!A \lif !B) \lif (\Box !A \lif \Box !B). \tag{\Ax{K}}
  \end{equation*}
\end{prop}

\begin{proof}
  Nous devons montrer que toutes les instances du schéma sont
  vraies en chaque monde de chaque modèle. Soient donc
  $\mModel{M} = \tuple{W,R,V}$ et $w \in W$ arbitraires.
  Pour établir qu'une implication est vraie en un monde, supposons
  son antécédent vrai et montrons que son conséquent l'est aussi.
  Supposons ici $\mSat{M}{\Box(!A \lif !B)}[w]$ et
  $\mSat{M}{\Box !A}[w]$; il faut montrer
  $\mSat{M}{\Box !B}[w]$. Soit $w'$ quelconque tel que $Rww'$.
  La première hypothèse donne $\mSat{M}{!A \lif !B}[w']$,
  et la seconde $\mSat{M}{!A}[w']$. Ainsi,
  $\mSat{M}{!B}[w']$. Comme $w'$ était quelconque,
  $\mSat{M}{\Box !B}[w]$.
\end{proof}

\begin{prop}\ollabel{prop:Dual-valid}
  Le schéma \Dual{} suivant est valide :
  \begin{equation*}
  \Diamond !A \liff \lnot\Box\lnot !A. \tag{\Dual}
  \end{equation*}
\end{prop}

\begin{proof}
  Exercice.
\end{proof}

\begin{prob}
  Démontrez la \olref[nml][syn][sch]{prop:Dual-valid}.
\end{prob}

\begin{prop}\ollabel{prop:soundMP}
  Si $!A$ et $!A \lif !B$ sont vraies en un monde d'un modèle,
  alors $!B$ y est vraie aussi. Par conséquent, les !!{formula}s
  valides sont closes par modus ponens.
\end{prop}

\begin{prop}\ollabel{prop:valid-instances}
  Une !!{formula}~$!A$ est valide si et seulement si toutes ses
  instances par substitution le sont. Autrement dit, un schéma
  est valide si et seulement si sa !!{formula} caractéristique
  l'est.
\end{prop}

\begin{proof}
  Le sens « si » est immédiat, puisque $!A$ est une instance
  d'elle-même par substitution.

  Pour le sens « seulement si », établissons d'abord le fait
  suivant. Soit $\mModel{M} = \tuple{W, R, V}$ un modèle modal,
  et soit $!B \ident
  \SSubst{!A}{\subst{!D_1}{p_1},\dots,\subst{!D_n}{p_n}}$ une
  instance de~$!A$ par substitution. Définissons
  $\mModel{M'} = \tuple{W, R, V'}$ en posant
  $V'(p_i) = \Setabs{w}{\mSat{M}{!D_i}[w]}$, et en attribuant
  arbitrairement des valeurs aux autres variables. Alors
  $\mSat{M}{!B}[w]$ si et seulement si
  $\mSat{M'}{!A}[w]$, pour tout~$w \in W$. (La démonstration
  est laissée en exercice.) Supposons maintenant que $!A$ soit
  valide, mais qu'une instance par substitution~$!B$ de~$!A$
  ne le soit pas.
  Il existe alors un modèle
  $\mModel{M} = \tuple{W, R, V}$ et un monde $w \in W$ tels que
  $\mSat/{M}{!B}[w]$. D'après le fait établi,
  $\mSat/{M'}{!A}[w]$ : $!A$ n'est donc pas valide,
  ce qui est contradictoire.
\end{proof}

\begin{prob}
  Démontrez le fait utilisé dans le sens « seulement si » de
  la preuve de la \olref[nml][syn][sch]{prop:valid-instances}.
  (Indication : procédez par induction sur~$!A$.)
\end{prob}

En revanche, un schéma n'est pas nécessairement vrai dans un
modèle lorsque sa formule caractéristique l'est. Le sens
« seulement si » demeure : si toute instance de $!A$ est vraie
dans~$\mModel{M}$, alors $!A$ elle-même est vraie
dans~$\mModel{M}$.
Mais $!A$ peut être vraie dans~$\mModel{M}$ alors qu'une
instance de~$!A$ est fausse dans un monde de~$\mModel{M}$.
Pour un contre-exemple simple, considérons~$p$ dans un modèle
ayant pour seul monde~$w$ et tel que $V(p) = \{w\}$. Ainsi,
$p$ est vraie en~$w$. Pourtant, $\lfalse$ est une instance
de~$p$ et elle n'est pas vraie en~$w$.

\begin{prob}
  Montrez qu'aucune des !!{formula}s suivantes n'est valide :
  \begin{enumerate}
    \item[\Ax{D}:] \quad $\Box p \lif \Diamond p$;
    \item[\Ax{T}:] \quad $\Box p \lif p$;
    \item[\Ax{B}:] \quad $p \lif \Box\Diamond p$;
    \item[\Ax{4}:] \quad $\Box p \lif \Box \Box p$;
    \item[\Ax{5}:] \quad $\Diamond p \lif \Box \Diamond
      p$.
  \end{enumerate}
\end{prob}

\begin{table}[t]
    \centering
    \begin{tabular}{| l || l |}
      \hline
      {\emph{Schémas valides}} & {\emph{Schémas non valides}} \\
      \hline\hline
      $\Box(!A \lif !B) \lif (\Diamond !A \lif \Diamond !B)$
      & $\Box (!A \lor !B) \lif (\Box !A \lor \Box !B)$ \\
      $\Diamond (!A \lif !B) \lif (\Box !A \lif \Diamond
      !B)$
      & $(\Diamond !A \land \Diamond !B) \lif \Diamond (!A
      \land !B)$\\
      $\Box (!A \land !B) \liff (\Box !A \land \Box !B)$
      & $!A \lif \Box !A$ \\
      $\Box !A \lif \Box (!B \lif !A)$
      & $\Box \Diamond !A \lif !B$ \\
      $\lnot \Diamond !A \lif \Box (!A \lif !B)$
      & $\Box \Box !A \lif \Box !A$ \\
      $\Diamond (!A \lor !B) \liff (\Diamond !A \lor
      \Diamond !B)$
      & $\Box \Diamond !A \lif \Diamond \Box !A$. \\
      \hline
    \end{tabular}
    \caption{Schémas valides et non valides.}
    \ollabel{tab:valid-invalidSchemas}
\end{table}

\begin{prob}%\ollabel{ex:in/validSchemas}
  Montrez que les schémas de la première colonne du
  \olref[nml][syn][sch]{tab:valid-invalidSchemas} sont valides,
  tandis que ceux de la seconde ne le sont pas.
\end{prob}

\begin{prob}
  Déterminez si les schémas suivants sont valides ou non :
  \begin{enumerate}
  \item $(\Diamond !A \lif \Box !B) \lif (\Box !A \lif \Box
    !B)$;
  \item $\Diamond(!A \lif !B) \lor \Box(!B \lif !A)$.
  \end{enumerate}
\end{prob}

\begin{prob}
  Pour chacun des schémas suivants, trouvez un modèle
  $\mModel{M}$ tel que toute instance de la !!{formula} soit
  vraie dans $\mModel{M}$ :
  \begin{enumerate}
  \item $p \lif \Diamond\Diamond p$;
  \item $\Diamond p \lif \Box p$.
  \end{enumerate}
\end{prob}

\end{document}
